\documentclass[11pt,a4paper]{article}

% ===== Packages =====
\usepackage[utf8]{inputenc}
\usepackage[T1]{fontenc}
\usepackage[english]{babel}
\usepackage{amsmath,amsthm,amssymb,amsfonts}
\usepackage{mathtools}
\usepackage{geometry}
\geometry{margin=2.5cm}
\usepackage{hyperref}
\hypersetup{
  colorlinks=true,
  linkcolor=blue!50!black,
  citecolor=green!50!black,
  urlcolor=blue!50!black
}
\usepackage{booktabs}
\usepackage{microtype}
\usepackage{xcolor}
\usepackage{listings}
\lstset{
  basicstyle=\ttfamily\small,
  breaklines=true,
  frame=single,
  framerule=0.3pt,
  framesep=4pt,
  columns=fullflexible,
  showstringspaces=false
}
\usepackage{algorithm}
\usepackage{algpseudocode}

% ===== Theorem environments =====
\theoremstyle{plain}
\newtheorem{theorem}{Theorem}[section]
\newtheorem{lemma}[theorem]{Lemma}
\newtheorem{proposition}[theorem]{Proposition}
\newtheorem{corollary}[theorem]{Corollary}

\theoremstyle{definition}
\newtheorem{definition}[theorem]{Definition}
\newtheorem{example}[theorem]{Example}
\newtheorem{remark}[theorem]{Remark}

\theoremstyle{plain}
\newtheorem*{conjecture*}{Conjecture}
\newtheorem{conjecture}[theorem]{Conjecture}

% ===== Macros =====
\newcommand{\NN}{\mathbb{N}}
\newcommand{\ZZ}{\mathbb{Z}}
\newcommand{\ZZp}[1]{\mathbb{Z}_{#1}}
\newcommand{\Branch}{\mathcal{B}}
\newcommand{\InvG}{G^{\leftarrow}}
\newcommand{\Calphax}{\mathcal{C}^{\leftarrow}_{S202}}
\newcommand{\aS}{a_{S202}}
\newcommand{\wS}{w_{S202}}
\newcommand{\defect}{\mathrm{D}}
\newcommand{\credit}{\mathrm{credit}}
\newcommand{\relevant}{\mathrm{relevant}}
\DeclareMathOperator{\InX}{InX}
\DeclareMathOperator{\InvEdge}{InvEdge}
\DeclareMathOperator{\modulus}{mod}

% ===== Title (double-blind: no authors) =====
\title{Bellman--Ford Closure Certificates for the
       Admissible Inverse Cylinder Graph of the Collatz Map\\
       \large{A formally verified barrier framework}}
\date{\today}
\author{}

\begin{document}
\maketitle

% Blind: no author block; section 0 acknowledges this.

\begin{abstract}
We introduce a new certificate framework for proving cylinder-accessibility
lower bounds in the admissible inverse tree of the Collatz dynamics.
The framework, called \emph{S216 Bellman--Ford optimistic closure},
addresses a subtle issue in earlier proposals: because some edges in
the inverse cylinder graph have negative weight, a naïve
``$d < T$'' relaxation cut fails to be a valid barrier criterion.
We replace it with the relevance predicate
\[
  \relevant(v, d) \iff d - (Q - j(v)) < T,
\]
which we show is the minimal saturation rule under which a finite
table over visited states proves $\Calphax(m, Q) \ge T$, the defect
lower bound on every admissible word landing in the $m$-th cylinder
of the published S202 fixed point.

The main theorem (Section~\ref{sec:closure-to-barrier}) is established
in two stages:
(i) a~\emph{suffix weight bound} certifying that the maximum future
descent from any vertex is at most $-(Q - j)$,
(ii) a~Bellman--Ford propagation argument showing that every prefix of
a hypothetical cheap path is \emph{relevant} and hence captured by
the closure invariant.

A complete formalization in the Lean~4 proof assistant is presented
in Section~\ref{sec:lean}, comprising
${\sim}1{,}600$ lines of new proof-script atop the public Mathlib
library, with $0$ \texttt{sorry}/\texttt{admit} axiomatic gaps and a
final theorem
\emph{Word $\Rightarrow$ barrier} chain proven end-to-end.
Section~\ref{sec:experiments} reports machine-verified certificates
for $(m, Q) \in \{(1,3), (1,5), (1,8), (1,10), (2,6)\}$, including
the case $Q = 10$ that previously exceeded the naïve budget.
\end{abstract}

\section{Introduction}\label{sec:intro}

The Collatz conjecture, dating back to 1937~\cite{lagarias2010,wirsching}, states that
the orbit of every positive integer under
\[
  T(n) =
    \begin{cases}
       n/2     & n \text{ even},\\
       (3n+1)/2 & n \text{ odd},
    \end{cases}
\]
eventually reaches $1$. After more than eight decades the conjecture
remains open; the strongest unconditional result is due to Tao~\cite{tao2022},
who proved that almost every Collatz orbit eventually attains an
\emph{almost bounded value}, building on prior probabilistic
arguments of Kontorovich--Sinai~\cite{kontorovich2002} and structural
analyses of Lagarias~\cite{lagarias1985}. The conjecture is
equivalent to the assertion that no positive integer admits an
infinite forward Collatz orbit disjoint from $\{1, 2\}$.
We work in the dual model of \emph{admissible inverse iteration}: in
this picture, words $u \in \{\mathtt{0},\mathtt{1}\}^*$ encode reverse
trajectories along the $\Psi_0,\Psi_1$ inverse branches, evaluated
from a starting integer $n_0 = 1$ via a state machine
\((n, A, q, B)\) carrying the translation identity
\[
   3^q\,n + B \;=\; 2^A.
\tag{Inv-Id}\label{eq:invid}
\]
The \emph{defect} of a word is $\defect(u) := A(u) - 2\,q(u)$.
Words with negative defect would yield a Collatz cycle off the trivial
$\{1,2,4\}$ cycle and refute the conjecture; conjecturally,
$\defect(u) \to \infty$ as $q(u) \to \infty$ along admissible
families.

A specific admissible word
$\wS = \mathtt{10100000001000010000000000}$ (the \emph{S202 word},
introduced in the companion preprint~\cite{S202})
is known to achieve a negative defect:
$(\defect, n, A, q) = (-1, 251, 43, 22)$ when evaluated from $n_0 = 1$.
Iterating $\wS$ asymptotically defines a 3-adic fixed point
$\aS = 1 + 3^{22} \in \ZZp{3}$, around which the cylinder family
\[
  C_m \;:=\; \{u : (n(u) \equiv \aS) \bmod 3^{22m+2}\}
\]
governs the asymptotic slope problem in the admissible cover.
The \emph{S202 alternative conjecture} asks whether
$C_m$ is accessible with $\defect = o(m)$ (subcritical slope), or
whether a structural obstruction blocks deep $C_m$ for $m \gg 1$.

The present work focuses on \emph{lower bounds}. We define
\[
  \Calphax(m, Q) :=
  \inf\bigl\{\, \defect(u) : u \in C_m \text{ and } q(u) \le Q\,\bigr\}.
\]
Establishing $\Calphax(m, Q) \ge T$ for concrete $(m, Q, T)$ is the
barrier program. Such inequalities translate, via the
concatenation lemma~\cite{S211}, into asymptotic slope statements:
if $\Calphax(m, Q) \ge T_m$ with $T_m / m \to c > 0$, then no
subcritical accessibility family can exist below slope $2 - c$.

\medskip

The main contribution of this paper is:

\begin{itemize}
\item A correct certificate format (Sections~\ref{sec:graph}--\ref{sec:closure-to-barrier})
      establishing $\Calphax(m, Q) \ge T$ from a finite, machine-checkable
      table of states; the format replaces the naïve cut by an
      \emph{optimistic} relevance rule.
\item A fully verified Lean~4 formalization (Section~\ref{sec:lean}) of
      the closure-to-barrier theorem together with the path-word
      correspondence (forward half) required to translate the abstract
      barrier into a defect lower bound for every admissible word
      landing in the cylinder.
\item Machine-verified certificates for the cases tabulated in
      Section~\ref{sec:experiments}, including a $352{,}716$-state
      certificate for $(m, Q) = (1, 10)$.
\end{itemize}

The artefact accompanying this submission contains the Python
generator (\texttt{s216\_colab.py}), the Lean~4 development
(\texttt{S216.lean}, \texttt{S212Forward.lean}, etc.), and the
auto-generated certificates.

\subsection*{Related work}\label{sec:related}
Formal verification of number-theoretic conjectures has grown rapidly
in the last decade, with notable milestones including the
formalization of the Kepler conjecture~\cite{flyspeck},
the prime number theorem~\cite{eberl-paulson}, and
Tao's recent informal-to-Lean projects~\cite{tao-blueprint}.
On the Collatz side, computer-assisted analyses go back to the early
1980s~\cite{lagarias1985}; the closest precedents to our approach
are the symbolic-dynamics framework of Conway~\cite{conway1972}
(reformulating Collatz as a halting problem for FRACTRAN), and the
3-adic analysis of Bernstein--Lagarias~\cite{bernstein-lagarias} which
underlies our cylinder framework.

Bellman--Ford on graphs with negative edges is classical in algorithm
design~\cite{cormen}; our novelty is the use of a problem-specific
\emph{credit potential} (Section~\ref{sec:closure-to-barrier}) that
preserves both decidability of the certificate and the formal
soundness of the closure-to-barrier implication under negative
weights.

\section{The Admissible Inverse Cylinder Graph}\label{sec:graph}

\subsection{Admissibility section}
Fix the admissibility section $\InX(n) :\iff n \bmod 9
\in \{1,2,4,5,7,8\}$, and the deterministic table
$\tau : \InX \to \{1,2,3,4\}$ given by
$\tau(1)=2,\ \tau(2)=1,\ \tau(4)=2,\ \tau(5)=3,\ \tau(7)=4,\ \tau(8)=1$.
The admissible state machine takes $(n,A,q,B) \in \NN^4$ to one of two
successor states according to the branch symbol.

\subsection{Inverse cylinder vertices}
Fix $m \in \NN$ and let $R := 22m + 2$ throughout. A vertex of the
inverse cylinder graph is a pair $(j, c) \in \NN^2$:
\[
  V_{m,Q} := \{(j, c) : 0 \le j \le Q\} \subset \NN^2,
\]
where $j$ counts the inverse \texttt{0}-edges already used and $c$
denotes the cylinder representative modulo $3^{R+j}$. The
\emph{canonical} version of $(j,c)$ satisfies $c < 3^{R+j}$.

The \emph{start} is
$s := (0, \aS \bmod 3^R)$;
a vertex $(j, c)$ is a \emph{goal} iff $c \equiv 1 \pmod{3^{R+j}}$.

\subsection{Edges}
There are two edge kinds (we abuse notation by reusing $v, v'$ for
canonical representatives):

\medskip
\noindent\textbf{Inverse-$\mathtt{1}$ edges.} For each admissible
$\alpha \in \{1,2,4,5,7,8\}$ with $\tau := \tau(\alpha)$, an edge
$v \to v'$ exists iff
\begin{enumerate}
  \item $j(v') = j(v)$,
  \item $v.c \equiv 2^{\tau} \cdot v'.c \pmod{3^{R+j}}$,
  \item $v'.c \equiv \alpha \pmod 9$,
\end{enumerate}
with weight $\omega = \tau$.

\medskip
\noindent\textbf{Inverse-$\mathtt{0}$ edges.}
For each admissible $\alpha$ and provided $j(v) < Q$,
\begin{enumerate}
  \item $j(v') = j(v) + 1$,
  \item $3v.c + 1 \equiv 2^{\tau(\alpha)} \cdot v'.c \pmod{3^{R+j(v')}}$,
  \item $v'.c \equiv \alpha \pmod 9$,
\end{enumerate}
with weight $\omega = \tau(\alpha) - 2$.

\medskip
\noindent The canonical $v'.c$ for each $\alpha$ is uniquely determined
by $v.c$:
\[
  v'.c
  \;=\;
  (2^\tau)^{-1} \cdot v.c
  \;\bmod\; 3^{R+j(v')},
\]
which we computed in practice using $\mathrm{invMod2}(M) := (M+1)/2$,
the closed-form inverse of $2$ modulo any odd $M$.

\subsection{Path weight and barrier}
A weighted directed path $p = v_0 \to v_1 \to \cdots \to v_k$ has
weight $w(p) := \sum_i \omega(v_{i-1}, v_i)$.
The cylinder barrier statement is
\[
  \Calphax(m, Q) \;:=\;
  \min_p w(p),
\]
the minimum over paths from $s$ to any goal vertex within the
$Q$-budget.

The structural correspondence (Theorem~\ref{thm:s212-forward}) states
that paths in this graph are in bijection with admissible words $u$
satisfying $q(u) \le Q$ and $n(u) \equiv \aS \pmod{3^R}$, with path
weight equal to $\defect(u)$. Hence the graph-theoretic barrier
becomes a defect lower bound on admissible words.

\section{S216: Bellman--Ford Closure with Negative Edges}\label{sec:closure-to-barrier}

\subsection{The negativity issue}
A naïve approach to prove $\Calphax(m, Q) \ge T$ enumerates all
states reachable from $s$ along paths of total weight $< T$ and
verifies that none is a goal. Equivalently, one builds a partial
dist function $\mathrm{dist} : V_{m,Q} \to \ZZ$ with closure
\[
  \forall v, \forall (v, v', \omega): \quad
  \mathrm{dist}(v') \le \mathrm{dist}(v) + \omega
\]
restricted to the region $\mathrm{dist}(v) < T$.

This is \emph{wrong} as soon as the graph admits edges of negative
weight. In our graph, the inverse-$\mathtt{0}$ edges with
$\alpha \in \{2, 8\}$ have $\tau = 1$ and therefore
$\omega = -1$. A path may temporarily accumulate weight above $T$ and
descend below it later. Restricting the saturation to
``$\mathrm{dist}(v) < T$'' misses precisely those prefixes that
exceed $T$ but contribute to a final goal of cost $< T$.

\subsection{The credit bound}
Negative edges are exactly the $j$-incrementing inverse-$\mathtt{0}$
steps. Hence from any vertex $v$, the number of negative-weight steps
remaining is bounded by $Q - j(v)$.

\begin{lemma}[Suffix weight bound]\label{lem:suffix}
For any path $v \to^* t$ confined to $V_{m,Q}$,
\[
  w(p) \;\ge\; -(Q - j(v)).
\]
\end{lemma}

\begin{proof}
Each edge has weight $\ge j(v_i) - j(v_{i+1})$
(positive edges have $\omega = \tau \ge 1$ and preserve $j$;
negative edges contribute $-1$ and increment $j$ by $1$).
Telescoping along the path,
$w(p) \ge j(v_0) - j(v_k) \ge -(Q - j(v_0))$,
using $j(v_k) \le Q$ on the bounded path.
\end{proof}

This is Lemma~6.1 in the supplementary Lean module
(\texttt{suffix\_weight\_bound}). Define the \emph{credit} of a
vertex as $\credit(Q, v) := Q - j(v)$.

\subsection{Optimistic relevance}
A state--cost pair $(v, d)$ is called \emph{relevant} (relative to
$Q, T$) iff
\[
  d - \credit(Q, v) \;<\; T.
\]
Otherwise the prefix $(v, d)$ provably cannot complete to a path of
weight $< T$ (Lemma~\ref{lem:suffix}), and may be pruned without loss.

\begin{lemma}[Prefix relevance]\label{lem:relevance}
If $p = s \to^* g$ is a path of total weight $w(p) < T$, then for
every prefix endpoint $v_i$ with accumulated cost $d_i$, the pair
$(v_i, d_i)$ is relevant.
\end{lemma}

\begin{proof}
$d_i + \text{(suffix weight from $v_i$ to $g$)} = w(p) < T$.
By Lemma~\ref{lem:suffix} applied to the suffix,
$d_i - \credit(Q, v_i) \le w(p) < T$.
\end{proof}

\subsection{Closure certificate}
\begin{definition}\label{def:cert}
An \emph{S216 closure certificate} for $(m, Q, T)$ is a partial
function $\mathrm{dist} : V_{m,Q} \to \ZZ \cup \{\bot\}$ together
with three invariants:
\begin{enumerate}
\item[(C1)] $\mathrm{dist}(s) = 0$.
\item[(C2)] No goal $g$ in the support of $\mathrm{dist}$ has
            $\mathrm{dist}(g) < T$.
\item[(C3)] For every vertex $v$ in the support of $\mathrm{dist}$
            with $\relevant(Q, T, v, \mathrm{dist}(v))$,
            and every outgoing edge $v \to v'$ of weight $\omega$
            with $j(v') \le Q$ and
            $\relevant(Q, T, v', \mathrm{dist}(v) + \omega)$,
            the value $\mathrm{dist}(v')$ is defined and satisfies
            $\mathrm{dist}(v') \le \mathrm{dist}(v) + \omega$.
\end{enumerate}
\end{definition}

\subsection{Main theorem}
\begin{theorem}[Closure-to-barrier]\label{thm:closure-to-barrier}
If an S216 closure certificate for $(m, Q, T)$ exists, then
$\Calphax(m, Q) \ge T$.
\end{theorem}

\begin{proof}
Suppose, for contradiction, that a path $p$ from $s$ to a goal $g$
has total weight $w(p) < T$. Let $v_0 = s, v_1, \dots, v_k = g$ be its
vertex sequence with prefix costs $d_i$ (so $d_0 = 0$ and $d_k = w(p)$).

We show by induction on $i$ that
$\mathrm{dist}(v_i)$ is defined and $\mathrm{dist}(v_i) \le d_i$.
The base case $i = 0$ holds by (C1).

For the inductive step, assume the property for $i$. By
Lemma~\ref{lem:relevance} the prefix $(v_i, d_i)$ is relevant. By
the inductive hypothesis $\mathrm{dist}(v_i) \le d_i$, so
$(v_i, \mathrm{dist}(v_i))$ is also relevant. Similarly
$(v_{i+1}, d_{i+1})$ is relevant, hence
$(v_{i+1}, \mathrm{dist}(v_i) + \omega_i)$ with
$\omega_i = w(v_i \to v_{i+1})$ is relevant (because
$\mathrm{dist}(v_i) + \omega_i \le d_{i+1}$).
By (C3), $\mathrm{dist}(v_{i+1})$ is defined and bounded above by
$\mathrm{dist}(v_i) + \omega_i \le d_i + \omega_i = d_{i+1}$.

Finally $\mathrm{dist}(g) \le d_k = w(p) < T$, contradicting (C2). \qed
\end{proof}

\section{From Barrier to Admissible Words}\label{sec:s212-forward}

The graph-theoretic barrier translates into a defect lower bound for
admissible words via the \emph{path-word correspondence}, the
``forward half'' of which we prove below.

\begin{theorem}[S212 forward]\label{thm:s212-forward}
Let $u$ be an admissible word with
$n(u) \equiv \aS \pmod{3^R}$.
Then there exists a weighted directed path in
$\InvG_{m,Q}$ from the start vertex $s$ to a canonical goal $g$ with
$j(g) = q(u)$ and weight $\defect(u)$.
Moreover, every intermediate vertex along the path is canonical.
\end{theorem}

\begin{proof}[Proof sketch]
By structural induction on $u$ from the front. The base case is
trivial (the empty word). For $u = b :: u'$, the inductive hypothesis
produces a path for $u'$ starting from the post-$b$ state; we
prepend an inverse-edge corresponding to the branch $b$, which is
exactly the predicate
$\InvEdge_{R}(v, v', \omega)$ for the explicit
$v, v', \omega$ supplied by the state-machine update.
Canonicality follows because every appended vertex is constructed as
$\bigl(j, x \bmod 3^{R+j}\bigr)$ for some $x \in \NN$.
The full proof is mechanized as \texttt{S212\_forward\_general} in
\texttt{S212Forward.lean}.
\end{proof}

\begin{corollary}\label{cor:full-chain}
Given $(m, Q, T)$ with an S216 closure certificate (Definition~\ref{def:cert})
whose start matches the inverse start, and given any admissible word $u$ with
$q(u) \le Q$ and $n(u) \equiv \aS \pmod{3^R}$, one has
$\defect(u) \ge T$.
\end{corollary}

\section{Formalization in Lean~4}\label{sec:lean}

The entire chain above is formalized in Lean~4~\cite{lean4} atop
Mathlib~\cite{mathlib}. The artefact consists of two new modules over
the existing \texttt{CollatzLean4} development:

\begin{itemize}
\item \texttt{S212Forward.lean}: forward path-word correspondence
      (Theorem~\ref{thm:s212-forward}) including canonicality of all
      path vertices.
\item \texttt{S216.lean}: definitions \texttt{credit}, \texttt{relevant};
      structures \texttt{S216Cert} and the Bool-decidable variant
      \texttt{S216Cert\_outgoing}; closure-to-barrier theorem; the
      structural lemmas \texttt{outgoingEdges\_sound},
      \texttt{outgoingEdges\_complete\_one}, and
      \texttt{outgoingEdges\_complete\_zero}.
\end{itemize}

\begin{lstlisting}
theorem closure_to_barrier_outgoing
    {R Q : Nat} {T : Int} (cert : S216Cert_outgoing R Q T)
    (hR : 1 ≤ R)
    (h_start_can : cert.start.c < 3 ^ (R + cert.start.j))
    {goal : InvVertex} {w : Int} {vs : List InvVertex}
    (h_path : WPath (InvEdge R) cert.start goal w vs)
    (h_g : InvVertex.IsGoal goal R)
    (h_start_bound : cert.start.j ≤ Q)
    (h_path_can : ∀ u ∈ vs, u.c < 3 ^ (R + u.j))
    (h_path_bound : ∀ u ∈ vs, u.j ≤ Q)
    (h_w_lt : w < T) :
    False
\end{lstlisting}

\noindent The full end-to-end barrier theorem is
\begin{lstlisting}
theorem S216_barrier_for_words_outgoing
    {m Q : Nat} {T : Int}
    (cert : S216Cert_outgoing (22 * m + 2) Q T)
    (h_start_match : cert.start = InvStart m)
    (h_start_can  : cert.start.c < 3 ^ (22 * m + 2 + cert.start.j))
    (u : Word)
    (h_q_bound : (evalWord u).q ≤ Q)
    (h_cyl : (evalWord u).n % (3 ^ (22 * m + 2))
              = aS202 % (3 ^ (22 * m + 2))) :
    T ≤ defect (evalWord u)
\end{lstlisting}

\noindent
The Lean development comprises ${\sim}1{,}600$ lines of new proof
script, built on top of Mathlib without any additional axiom and with
$0$ \texttt{sorry} placeholders.

The \emph{Bool-decidable} variant \texttt{S216Cert\_outgoing} replaces
the \texttt{InvEdge} quantifier of Definition~\ref{def:cert} by
iteration over the explicit \texttt{outgoingEdges} enumerator. The
soundness and canonical completeness lemmas certify that this
enumeration is equivalent to the abstract \texttt{InvEdge} predicate
for the canonical paths produced by Theorem~\ref{thm:s212-forward}.

\section{Algorithm and Experimental Verification}\label{sec:experiments}

\subsection{The certificate-generation algorithm}
Algorithm~\ref{alg:bf} produces an S216 certificate (or a refutation,
i.e.\ a cheap path to a goal) by an SPFA-style label-correcting
relaxation, pruning any label that fails the optimistic relevance
predicate. Termination on a finite candidate is guaranteed since the
relevant region is finite (bounded by $Q$ and by $T + \credit(Q,v)$).

\begin{algorithm}[h]
\caption{S216 BF closure (relevance-pruned label correction)}\label{alg:bf}
\begin{algorithmic}[1]
\State $\mathrm{dist}[s] \gets 0$;\ \ $\mathit{queue} \gets [\,s\,]$
\While{$\mathit{queue}$ non-empty}
  \State $v \gets$ pop front of $\mathit{queue}$;\ $d \gets \mathrm{dist}[v]$
  \If{$d - (Q - j(v)) \ge T$} \textbf{continue} \Comment{prune non-relevant}
  \EndIf
  \If{$v$ is a goal and $d < T$} \Return $(\textsc{useful\_path}, v, d)$
  \EndIf
  \For{$(v', \omega) \in \mathrm{outgoingEdges}(v)$}
    \State $d' \gets d + \omega$
    \If{$v'$ is goal and $d' < T$} \Return $(\textsc{useful\_path}, v', d')$
    \EndIf
    \If{$d' - (Q - j(v')) \ge T$} \textbf{continue} \Comment{non-relevant target}
    \EndIf
    \If{$\mathrm{dist}[v']$ undefined or $d' < \mathrm{dist}[v']$}
      \State $\mathrm{dist}[v'] \gets d'$;\ \ enqueue $v'$
    \EndIf
  \EndFor
\EndWhile
\State \Return $(\textsc{barrier}, \mathrm{dist})$
\end{algorithmic}
\end{algorithm}

The state space size scales roughly as $O(6^Q)$ in the worst case
(six admissible $\alpha$ choices per zero-step); empirically the
relevant region grows much more slowly because most successor
candidates fail $\alpha$-admissibility or relevance.

\subsection{Experimental certificates}
Table~\ref{tab:certs} reports certificate sizes and verification
status for several $(m, Q)$. The Python script
\texttt{s216\_colab.py} implements Algorithm~\ref{alg:bf} together
with an independent verifier that re-checks (C1)--(C3) against a
freshly loaded certificate JSON.

\begin{table}[h]
\centering
\caption{S216 BF closure certificates. ``Naïve'' is the legacy
$d < T$ cut (insufficient for a formal barrier); ``S216'' is the
optimistic cut $d - (Q-j) < T$. All S216 certificates verify
independently. The case $(1, 10)$ exceeded the local budget under the
S202 engine but completes under the Colab pipeline.}
\label{tab:certs}
\begin{tabular}{rrrrrr}
\toprule
$m$ & $Q$ & $T$ & Naïve states & S216 states & verifier \\
\midrule
1 & 3  & 1 & 22                 & 35           & ok \\
1 & 5  & 1 & 196                & 462          & ok \\
1 & 8  & 1 & 6{,}917            & 24{,}310     & ok \\
1 & 10 & 1 & $>\!200{,}000^{*}$ & 352{,}716    & ok \\
2 & 6  & 2 & 1{,}429            & 3{,}003      & ok \\
\bottomrule
\end{tabular}
\end{table}

\noindent ($*$) The naïve cut produces an incomplete table within the
$2 \times 10^5$ budget. Even with a higher budget, the resulting
table would not constitute a valid barrier under S216 since negative
prefixes are pruned.

\subsection{Reproducibility}
The full artefact builds under Lean~4.30.0-rc2 with the Mathlib
revision pinned by the accompanying \texttt{lake-manifest.json}. The
Python pipeline runs on Python~3.10 or higher with no external
dependencies beyond the standard library. A representative invocation:
\begin{lstlisting}
$ python tools/s216_colab.py --m 1 --Q 10 --budget 5000000
=== S216 engine: m=1 Q=10 T=1 R=24 ===
[s216] states=352716 expanded=352716 edges=727856 t=2.5s
VERIFICATION: ok
Conclusion: C^<-_{S202}(m=1, Q=10) >= 1
$ lake build CollatzLean4
Build completed successfully (8415 jobs).
\end{lstlisting}
A fresh build of the Lean development completes in under two minutes
on commodity hardware. All five certificates in Table~\ref{tab:certs}
are reproduced in under a minute aggregate.

\section{Open Problems}\label{sec:open}

\subsection*{The reverse half of the path--word correspondence}
We have proven the forward direction
(Theorem~\ref{thm:s212-forward}); the reverse direction states that
every inverse-graph path with $j(g) = q$ corresponds to at least one
admissible word with $\defect = w$. Combining both halves would yield
the full bijection (the conjectural \texttt{S212\_path\_word\_correspondence}
in the Lean development).

\subsection*{Analytic potentials}
The S216 certificate is an explicit lookup table over visited
canonical states, of size growing roughly as a polynomial in $Q$. A
much cleaner \emph{analytic potential} would be a closed-form
function $\Phi : V_{m,Q} \to \ZZ$ with the local potential
inequality
$\Phi(v) \le \omega + \Phi(v')$ along every edge,
combined with start/goal boundary conditions. The S214
conjecture~\cite{S214} (formalized in the supplementary Lean module
\texttt{Potential.lean}) proposes a mixed analytic potential of the form
\[
  \Phi(j, c) = \lambda\,\delta_+(c, R+j)
             + \gamma\,\rho_-(c, R+j)
             - \mu(Q - j)
             + \psi(c \bmod 3^h),
\]
where $\delta_+,\rho_-$ are 3-adic distance functions to $1$ and $-1$.
Finding $\lambda,\gamma,\mu,\psi$ satisfying the edge inequality is
the key analytic problem.

\subsection*{Scaling laws of $\Calphax(m, Q)$}
The most pressing open question is the asymptotic of
$\Calphax(m, Q)$ as $Q \to \infty$ for fixed $m$. The S202
alternative conjecture amounts to whether $\Calphax(m, Q)$ remains
unbounded (slope barrier) or stabilizes (subcritical accessibility).

% =====================================================================
% S214-Core companion section (3-adic accessibility potentials).
% Defined in paper/s214_core_section.tex; also compiles standalone via
% paper/s214_core.tex.
% =====================================================================
% =====================================================================
% S214-Core: 3-adic accessibility potentials
% Standalone section fragment. \input from a host document that defines
% the lemma/proposition/conjecture/proof environments and loads amsmath.
% =====================================================================

\section{S214-Core: 3-adic Accessibility Potentials}
\label{sec:S214-core}

The purpose of this section is to isolate the structural mechanism behind
the accessibility barrier for the S202 block. The block S202 shows that
local subcriticality in the admissible inverse system is possible:
\[
A(w_{202})-2q(w_{202})=-1.
\]
Thus the naive inequality \(A\ge 2q\) is false. The correct question is no
longer whether subcritical blocks exist, but whether such blocks are
accessible from the root with enough gain to overcome their entrance cost.

The guiding principle is:

\[
\boxed{
\text{Subcritical }3\text{-adic blocks may exist, but they may be protected
by accessibility barriers.}
}
\]

The first test case is the S202 block.

% ---------------------------------------------------------------------
\subsection{The S202 fixed cylinder}
\label{subsec:S202-cylinder}

Let
\[
w_{202}=10100000001000010000000000,
\]
with certified data
\[
A_{202}=43,\qquad q_{202}=22,\qquad D_{202}=A_{202}-2q_{202}=-1.
\]
The associated affine action is
\[
F_{202}(x)=\frac{2^{43}x-B_{202}}{3^{22}},
\]
where
\[
B_{202}=919447060349.
\]
Its formal fixed point in \(\mathbb Z_3\) is
\[
a_{202}
=
\frac{B_{202}}{2^{43}-3^{22}}.
\]

A direct computation gives
\[
a_{202}\equiv 1+3^{22}\pmod{3^{24}}.
\]
Therefore
\[
v_3(a_{202}-1)=22.
\]

For \(m\ge1\), define
\[
R_m=22m+2
\]
and the S202 target cylinder
\[
C_m
=
a_{202}+3^{R_m}\mathbb Z_3.
\]

The accessibility cost is
\[
\mathcal C^{\leftarrow}_{S202}(m,Q)
=
\inf
\left\{
D(u):
q(u)\le Q,\
n(u)\equiv a_{202}\pmod{3^{R_m}}
\right\}.
\]

Since concatenating \(m\) copies of \(w_{202}\) decreases the defect by \(m\),
we have
\[
D(u w_{202}^{m})
=
D(u)-m.
\]
Hence, if
\[
\mathcal C^{\leftarrow}_{S202}(m,Q)<m,
\]
then there exists an accessible concatenation with negative defect:
\[
D(u w_{202}^{m})<0.
\]

Conversely, if
\[
\mathcal C^{\leftarrow}_{S202}(m,Q)\ge m,
\]
then S202 cannot produce a net subcritical concatenation under the budget
\(q(u)\le Q\).

% ---------------------------------------------------------------------
\subsection{Root and negative-run proximities}
\label{subsec:rho-plus-minus}

Let \([c]_r\) denote the \(3\)-adic cylinder
\[
[c]_r=c+3^r\mathbb Z_3.
\]

We define the root proximity
\[
\rho_+(c,r)
=
\min\{v_3(c-1),r\},
\]
and the root deficit
\[
\delta_+(c,r)
=
r-\rho_+(c,r).
\]

Thus
\[
\delta_+(c,r)=0
\]
if and only if
\[
[c]_r
\]
contains the root \(1\).

We also define the negative-run proximity
\[
\rho_-(c,r)
=
\min\{v_3(c+1),r\}.
\]

The quantity \(\rho_-(c,r)\) measures proximity to the \(3\)-adic point
\(-1\), which controls the possibility of long negative defect runs.

% ---------------------------------------------------------------------
\subsection{Separation between the root world and the negative-run world}
\label{subsec:root-negative-separation}

\begin{lemma}[Root/negative separation]
\label{lem:root-negative-separation}
For every integer \(c\) and every \(r\ge1\),
\[
\min\bigl(\rho_+(c,r),\rho_-(c,r)\bigr)=0.
\]
Equivalently, a \(3\)-adic cylinder cannot be simultaneously close to
\(1\) and to \(-1\).
\end{lemma}

\begin{proof}
Suppose first that \(\rho_+(c,r)>0\). Then
\[
v_3(c-1)>0,
\]
so
\[
c\equiv1\pmod3.
\]
Therefore
\[
c+1\equiv2\pmod3,
\]
and hence
\[
v_3(c+1)=0.
\]
Thus
\[
\rho_-(c,r)=0.
\]

Similarly, if \(\rho_-(c,r)>0\), then
\[
c\equiv -1\pmod3,
\]
so
\[
c-1\equiv -2\not\equiv0\pmod3.
\]
Hence
\[
v_3(c-1)=0,
\]
and therefore
\[
\rho_+(c,r)=0.
\]

Thus the two truncated valuations cannot both be positive.
\end{proof}

This elementary separation is one of the central structural facts of the
program. It means that paths attempting to exploit negative defect gains
must move into the \(-1\)-world, whereas paths attempting to terminate at
the root must move into the \(1\)-world. These two regions are disjoint
already modulo \(3\).

% ---------------------------------------------------------------------
\subsection{Negative defect steps and proximity to \(-1\)}
\label{subsec:negative-runs}

Recall that the defect is
\[
D=A-2q.
\]
Under a branch \(1\),
\[
D\mapsto D+\tau(n),
\]
while under a branch \(0\),
\[
D\mapsto D+\tau(n)-2.
\]

The only negative defect steps occur when
\[
\tau(n)=1
\]
and the branch is \(0\). This happens precisely for
\[
n\equiv2,8\pmod9.
\]

In that case the map is
\[
L(n)=\frac{2n-1}{3}.
\]

\begin{lemma}[Iterated negative branch]
\label{lem:iterated-negative-branch}
For every \(k\ge0\),
\[
L^k(n)
=
\frac{2^k n+2^k-3^k}{3^k}.
\]
Equivalently,
\[
L^k(n)+1
=
\frac{2^k(n+1)}{3^k}.
\]
\end{lemma}

\begin{proof}
For \(k=0\), the formula is immediate. Assume it holds for \(k\). Then
\[
L^{k+1}(n)
=
L(L^k(n))
=
\frac{2L^k(n)-1}{3}.
\]
Using the induction hypothesis,
\[
2L^k(n)-1
=
2\cdot
\frac{2^k n+2^k-3^k}{3^k}
-1.
\]
Thus
\[
2L^k(n)-1
=
\frac{2^{k+1}n+2^{k+1}-2\cdot3^k-3^k}{3^k}
=
\frac{2^{k+1}n+2^{k+1}-3^{k+1}}{3^k}.
\]
Dividing by \(3\), we obtain
\[
L^{k+1}(n)
=
\frac{2^{k+1}n+2^{k+1}-3^{k+1}}{3^{k+1}}.
\]
The equivalent formula follows by adding \(1\):
\[
L^k(n)+1
=
\frac{2^k n+2^k}{3^k}
=
\frac{2^k(n+1)}{3^k}.
\]
\end{proof}

\begin{lemma}[Negative runs force proximity to \(-1\)]
\label{lem:negative-runs-force-minus-one}
If \(k\) consecutive negative steps are defined starting from \(n\), then
\[
n\equiv -1\pmod{3^k}.
\]
Equivalently,
\[
v_3(n+1)\ge k.
\]
\end{lemma}

\begin{proof}
If \(k\) consecutive applications of
\[
L(n)=\frac{2n-1}{3}
\]
are defined over the integers, then \(L^k(n)\in\mathbb Z\). By
Lemma~\ref{lem:iterated-negative-branch},
\[
L^k(n)+1
=
\frac{2^k(n+1)}{3^k}.
\]
Since \(2^k\) is invertible modulo powers of \(3\), integrality forces
\[
3^k\mid n+1.
\]
Hence
\[
n\equiv -1\pmod{3^k}.
\]
\end{proof}

The converse requires the fine admissibility condition: integrality alone
is not sufficient to guarantee that all intermediate states remain in the
negative classes \(2,8\pmod9\). More precisely, a genuine negative run of
length \(k\) corresponds to two of the three classes above
\[
-1\pmod{3^k}.
\]
Equivalently, one needs
\[
n\equiv -1+3^k m\pmod{3^{k+1}},
\]
with
\[
m\in\{0,1,2\},
\qquad
m\not\equiv 2^k\pmod3.
\]

This refinement is important computationally, but the coarse implication
\[
\text{negative run of length }k
\Longrightarrow
v_3(n+1)\ge k
\]
is already enough to justify the appearance of \(\rho_-\) in the potential.

% ---------------------------------------------------------------------
\subsection{The S202 root-distance deficit}
\label{subsec:S202-root-distance-deficit}

\begin{lemma}[S202 root-distance deficit]
\label{lem:S202-root-deficit}
For the S202 fixed point and
\[
R_m=22m+2,
\]
one has
\[
\rho_+(a_{202},R_m)=22,
\]
and therefore
\[
\delta_+(a_{202},R_m)
=
R_m-22
=
22m-20.
\]
\end{lemma}

\begin{proof}
The certified computation gives
\[
a_{202}\equiv1+3^{22}\pmod{3^{24}}.
\]
Thus
\[
a_{202}-1
=
3^{22}u
\]
for some \(3\)-adic unit \(u\), since the coefficient of \(3^{22}\) is
nonzero modulo \(3\). Hence
\[
v_3(a_{202}-1)=22.
\]
For \(m\ge1\),
\[
R_m=22m+2\ge24
\]
except for the minimal boundary where the same valuation conclusion remains
valid by truncation. Therefore
\[
\rho_+(a_{202},R_m)
=
\min\{22,R_m\}
=
22.
\]
Consequently,
\[
\delta_+(a_{202},R_m)
=
R_m-\rho_+(a_{202},R_m)
=
22m+2-22
=
22m-20.
\]
\end{proof}

This lemma is the first quantitative indication of an accessibility
barrier. Although \(a_{202}\) is close to the root \(1\) modulo \(3^{22}\),
the cylinder \(C_m\) requires agreement modulo
\[
3^{22m+2}.
\]
Thus the missing root precision is
\[
22m-20.
\]

The S202 block gives only \(m\) units of defect gain after \(m\)
concatenations. Hence the accessibility problem asks whether the system can
correct approximately \(22m\) missing \(3\)-adic digits while paying less
than \(m\) units of defect.

% ---------------------------------------------------------------------
\subsection{Candidate accessibility potentials}
\label{subsec:candidate-potentials}

The previous lemmas suggest that any useful path toward the S202 cylinder
must balance two competing mechanisms:

\[
\boxed{
\text{approaching the root }1
}
\]
and
\[
\boxed{
\text{exploiting negative runs near }-1.
}
\]

This motivates potentials of the form
\[
\Phi(j,c)
=
\lambda\,\delta_+(c,r)
+
\gamma\,\rho_-(c,r)
-
\mu(Q-j)
+
\psi(c\bmod 3^h),
\]
where
\[
r=R_m+j.
\]

Here:

\begin{itemize}
    \item \(j\) is the number of inverse zero-steps already used;
    \item \(Q-j\) is the remaining zero budget;
    \item \(\delta_+(c,r)\) measures distance from the root \(1\);
    \item \(\rho_-(c,r)\) measures proximity to the negative-run region
          around \(-1\);
    \item \(\psi(c\bmod3^h)\) is a finite local correction.
\end{itemize}

The role of the term
\[
-\mu(Q-j)
\]
is to account for the maximum possible future defect decrease, since each
remaining zero can reduce the defect by at most \(1\).

The role of
\[
\gamma\rho_-(c,r)
\]
is to penalize excessive reliance on negative runs: long negative runs
require proximity to \(-1\), but by
Lemma~\ref{lem:root-negative-separation}, this is incompatible with
proximity to the root \(1\).

The finite correction
\[
\psi(c\bmod3^h)
\]
is necessary because the evolution of \(\rho_+\) and \(\rho_-\) under the
inverse branches is not determined solely by their current values. Local
residue data modulo \(3^h\) is required.

% ---------------------------------------------------------------------
\subsection{Dual certificate condition}
\label{subsec:dual-certificate-condition}

Let \(G_{m,Q}\) be the inverse cylinder graph associated with the S202
target cylinder and zero budget \(Q\). A vertex is written
\[
v=(j,c),
\]
and each edge
\[
v\to v'
\]
has weight
\[
\omega(v,v')\in\{-1,0,1,2,3,4\}.
\]

A potential \(\Phi\) certifies a lower bound if:

\[
\Phi(g)\le0
\]
for every goal vertex \(g\), and
\[
\Phi(v)\le \omega(v,v')+\Phi(v')
\]
for every edge \(v\to v'\).

\begin{proposition}[Dual barrier certificate]
\label{prop:dual-barrier-certificate}
If a potential \(\Phi\) satisfies
\[
\Phi(g)\le0
\]
for every goal vertex \(g\), and
\[
\Phi(v)\le \omega(v,v')+\Phi(v')
\]
for every edge \(v\to v'\), then every path from \(s\) to a goal has weight
at least
\[
\Phi(s).
\]
In particular, if
\[
\Phi(s)\ge T,
\]
then
\[
\mathcal C^{\leftarrow}_{S202}(m,Q)\ge T.
\]
\end{proposition}

\begin{proof}
Let
\[
s=v_0\to v_1\to\cdots\to v_k=g
\]
be a path from the start vertex \(s\) to a goal vertex \(g\). Applying the
edge inequality to every edge gives
\[
\Phi(v_i)
\le
\omega(v_i,v_{i+1})+\Phi(v_{i+1}).
\]
Summing telescopically,
\[
\Phi(s)
\le
\sum_{i=0}^{k-1}\omega(v_i,v_{i+1})+\Phi(g).
\]
Since \(g\) is a goal and \(\Phi(g)\le0\), we get
\[
\Phi(s)
\le
\sum_{i=0}^{k-1}\omega(v_i,v_{i+1}).
\]
Thus every path from \(s\) to a goal has weight at least \(\Phi(s)\).
\end{proof}

This proposition is the abstract form of the S216 closure-to-barrier
certificate. The Bellman--Ford closure certificate is a finite,
machine-checkable way of constructing such a potential or, equivalently,
proving that no path of weight \(<T\) exists.

% ---------------------------------------------------------------------
\subsection{S202 accessibility barrier conjecture}
\label{subsec:S202-barrier-conjecture}

The computations for S202 suggest the following conjecture.

\begin{conjecture}[S202 accessibility barrier]
\label{conj:S202-accessibility-barrier}
For the S202 block, there exists a constant \(K>0\) such that, for all
sufficiently large \(m\),
\[
\mathcal C^{\leftarrow}_{S202}(m,Km)\ge m.
\]
Equivalently, no access path using at most \(Km\) inverse zero-steps can
enter the cylinder
\[
a_{202}+3^{22m+2}\mathbb Z_3
\]
with defect smaller than the \(m\) units of defect gained by
\(w_{202}^{m}\).
\end{conjecture}

A stronger version would assert that there exists \(\varepsilon>0\) such
that
\[
\mathcal C^{\leftarrow}_{S202}(m,Km)
\ge
(1+\varepsilon)m
\]
for all sufficiently large \(m\).

The meaning of the conjecture is:

\[
\boxed{
\text{S202 is a subcritical }3\text{-adic block, but it is not cheaply
accessible from the root.}
}
\]

% ---------------------------------------------------------------------
\subsection{Universal subcritical accessibility barrier}
\label{subsec:universal-subcritical-barrier}

The S202 block should be regarded as the first test case of a broader
principle.

Let \(w\) be any admissible block with data
\[
A(w),\qquad q(w),\qquad B(w),
\]
and defect
\[
D(w)=A(w)-2q(w).
\]

Suppose
\[
D(w)<0.
\]
Write
\[
\Delta(w)=-D(w)>0.
\]

If \(w\) has a \(3\)-adic fixed point
\[
a_w=
\frac{B(w)}{2^{A(w)}-3^{q(w)}}\in\mathbb Z_3,
\]
then the natural analogue of the S202 barrier is
\[
\mathcal C^{\leftarrow}_{w}(m,Q)
\ge
m\Delta(w).
\]

\begin{conjecture}[Universal subcritical accessibility barrier]
\label{conj:universal-subcritical-barrier}
For every admissible subcritical block \(w\) with \(D(w)<0\), there exists a
linear budget \(Q_m=O(m)\) such that, for all sufficiently large \(m\),
\[
\mathcal C^{\leftarrow}_{w}(m,Q_m)
\ge
-mD(w).
\]
Equivalently, the cost of accessing the \(m\)-th fixed cylinder of a
subcritical block is at least the total defect gain obtained by repeating
that block \(m\) times.
\end{conjecture}

This conjecture expresses the general principle:

\[
\boxed{
\text{local subcriticality does not imply globally accessible
subcriticality.}
}
\]

% ---------------------------------------------------------------------
\subsection{Relation to the Collatz problem}
\label{subsec:relation-to-collatz}

The S202 accessibility barrier, even if proved, does not by itself prove
the Collatz conjecture. It only proves that one specific subcritical
\(3\)-adic block cannot be exploited cheaply from the root.

To obtain a route toward Collatz, one needs an additional reduction
principle.

\begin{conjecture}[Subcritical reduction of counterexamples]
\label{conj:subcritical-reduction}
If the Collatz conjecture fails, then there exists an accessible family of
admissible inverse words whose asymptotic defect becomes negative after
entering a recurrent subcritical block or family of blocks.
\end{conjecture}

In schematic form, the desired strategy is:

\[
\boxed{
\text{Counterexample}
\Longrightarrow
\text{accessible subcritical family}
}
\]
and
\[
\boxed{
\text{universal accessibility barrier}
\Longrightarrow
\text{no accessible subcritical family}.
}
\]

Together these would imply the nonexistence of a counterexample.

Thus the long-term program is:

\[
\boxed{
\text{Reduction of counterexamples}
+
\text{universal subcritical accessibility barrier}
\Longrightarrow
\text{Collatz}.
}
\]

At the current stage, S202 provides the first nontrivial test case for this
program, and S214 provides the potential-theoretic language needed to move
from finite certificates to structural barriers.

% ---------------------------------------------------------------------
\subsection{Immediate formal milestones}
\label{subsec:S214-formal-milestones}

The immediate formal targets are:

\begin{enumerate}
    \item Complete the S216 closure-to-barrier bridge in Lean.
    \item Obtain the first formal theorem:
    \[
    \mathcal C^{\leftarrow}_{S202}(1,3)\ge1.
    \]
    \item Scale the certificate generation to:
    \[
    (m,Q,T)=(1,5,1),\ (1,8,1),\ (1,10,1),\ (2,6,2).
    \]
    \item Extract empirical profiles of the cheap region using:
    \[
    \rho_+,\qquad \rho_-,\qquad j,\qquad D.
    \]
    \item Use these profiles to search for explicit finite-depth potentials
    of the form
    \[
    \Phi(j,c)
    =
    \lambda\delta_+(c,r)
    +
    \gamma\rho_-(c,r)
    -
    \mu(Q-j)
    +
    \psi(c\bmod3^h).
    \]
\end{enumerate}

The first fully verified Lean theorem
\[
\mathcal C^{\leftarrow}_{S202}(1,3)\ge1
\]
will constitute the first formal accessibility barrier in the S202 chain.


\section{Conclusion}

We have presented an explicitly machine-checkable barrier framework
for the cylinder-accessibility problem in the admissible inverse
Collatz tree, fully formalized in Lean~4 with no axiomatic gaps.
The key analytical insight is the credit-based relevance rule
(Definition~\ref{def:cert}, item (C3)), which corrects a subtle
prune-bug in earlier proposals and renders the certificate suitable
both for machine verification and for formal lifting.

The artefact accompanying this submission provides:
(i) the Python certificate generator and independent verifier,
(ii) a complete Lean~4 formalization,
(iii) sample certificates for $(m, Q) \in \{(1,3),\dots,(2,6)\}$.

\bibliographystyle{plain}
\begin{thebibliography}{99}

% --- Collatz literature ---
\bibitem{lagarias1985}
J.~C. Lagarias.
\newblock The 3$x$ + 1 problem and its generalizations.
\newblock {\em American Mathematical Monthly}, 92(1):3--23, 1985.

\bibitem{lagarias2010}
J.~C. Lagarias (editor).
\newblock {\em The Ultimate Challenge: The 3$x$+1 Problem}.
\newblock American Mathematical Society, 2010.

\bibitem{wirsching}
G.~J. Wirsching.
\newblock {\em The Dynamical System Generated by the 3$n$+1 Function}.
\newblock Lecture Notes in Mathematics 1681, Springer, 1998.

\bibitem{kontorovich2002}
A.~V. Kontorovich and Ya.~G. Sinai.
\newblock Structure theorem for $(d,g,h)$-maps.
\newblock {\em Bulletin of the Brazilian Mathematical Society},
33(2):213--224, 2002.

\bibitem{tao2022}
T.~Tao.
\newblock Almost all orbits of the Collatz map attain almost bounded
values.
\newblock {\em Forum of Mathematics, Pi}, 10(e12), 2022.

\bibitem{conway1972}
J.~H. Conway.
\newblock Unpredictable iterations.
\newblock In {\em Proceedings of the 1972 Number Theory Conference},
49--52. University of Colorado, 1972.

\bibitem{bernstein-lagarias}
D.~J. Bernstein and J.~C. Lagarias.
\newblock The 3$x$+1 conjugacy map.
\newblock {\em Canadian Journal of Mathematics}, 48(6):1154--1169, 1996.

% --- Formal verification ---
\bibitem{flyspeck}
T.~C. Hales \emph{et al}.
\newblock A formal proof of the Kepler conjecture.
\newblock {\em Forum of Mathematics, Pi}, 5(e2), 2017.

\bibitem{eberl-paulson}
M.~Eberl, L.~C. Paulson.
\newblock A formal proof of the prime number theorem in Isabelle/HOL.
\newblock {\em Archive of Formal Proofs}, 2018.

\bibitem{tao-blueprint}
T.~Tao \emph{et al}.
\newblock Equational theories project: a Lean blueprint methodology
for collaborative formal mathematics.
\newblock Online, \url{https://teorth.github.io/equational_theories},
2024--2026.

\bibitem{lean4}
L.~de~Moura and S.~Ullrich.
\newblock The Lean~4 theorem prover and programming language.
\newblock In {\em CADE-28}, 2021.

\bibitem{mathlib}
The mathlib Community.
\newblock The Lean Mathematical Library.
\newblock In {\em CPP 2020}, 367--381, 2020.
\newblock \url{https://leanprover-community.github.io}.

% --- Algorithms ---
\bibitem{cormen}
T.~H. Cormen, C.~E. Leiserson, R.~L. Rivest, and C.~Stein.
\newblock {\em Introduction to Algorithms}, 3rd edition.
\newblock MIT Press, 2009.

% --- Within-program references (self-citations, anonymised) ---
\bibitem{S202}
\emph{[Author manuscript]}.
S202 tree: admissible inverse iteration around the
$\protect\aS = 1+3^{22}$ fixed point.
Companion preprint, available in artefact supplement, 2026.

\bibitem{S211}
\emph{[Author manuscript]}.
S211 concatenation lemma: subcritical accessibility
implies subcritical slope.
Companion preprint, available in artefact supplement, 2026.

\bibitem{S214}
\emph{[Author manuscript]}.
S214: local 3-adic potentials in the admissible
inverse cylinder graph.
Companion preprint, available in artefact supplement, 2026.

\end{thebibliography}

\end{document}
